\begingroup
\setlength{\abovedisplayskip}{4pt plus 1pt minus 1pt}
\setlength{\belowdisplayskip}{4pt plus 1pt minus 1pt}
\setlength{\abovedisplayshortskip}{2pt plus 1pt}
\setlength{\belowdisplayshortskip}{2pt plus 1pt}
\setlength{\parskip}{0pt}
\section{\HMTLang{Mémoire bilatérale, holonomie relative et restitution orientée}{Mémoire bilatérale, holonomie relative et récupération orientée}}
\label{delta22:xi:holonomia}
\HMTLang{La conservation bilatérale acquiert un contenu opérationnel lorsque l’on conserve, en plus de l’état visible, le canal qui enregistre la composition. La sélection exceptionnelle et temporelle de \(K\), développée précédemment, détermine le registre ; ses lecteurs arithmétiques et incidenciels transmettent des relations distinctes du même état enrichi. On explicite ici comment un ordre d’opérations laisse une réponse récupérable même lorsque leurs multiplicités coïncident et que la phase revient.}{La conservation bilatérale acquiert un contenu opérationnel lorsque le canal enregistrant la composition est retenu aux côtés de l’état visible. La sélection exceptionnelle et temporelle de \(K\), développée ci-dessus, détermine le registre ; ses lecteurs arithmétiques et d’incidence transmettent des relations distinctes du même état enrichi. Nous explicitons ici comment un ordre d’opérations laisse une réponse récupérable même lorsque les multiplicités concordent et que la phase revient.}
\HMTLang{L’antécédent est la génération APP--TRIT--TPK, avec ses deux projections, le prolongement complet et l’orientation dodécaphasique. L’holonomie nonadique \(\Gamma_9\) appartient à cette dynamique enrichie. Le commutateur relatif construit dans ce paragraphe correspond à une réalisation opératoire spécifiée ; son domaine, sa forme conservée et son lecteur restent distingués.}{L’antécédent est la génération APP--TRIT--TPK, avec les deux projections, le prolongement complet et l’orientation dodécaphasique. L’holonomie nonadique \(\Gamma_9\) appartient à cette dynamique enrichie. Le commutateur relatif construit dans cette section appartient à une réalisation opératorielle spécifiée ; son domaine, sa forme conservée et son lecteur restent distingués.}

\subsection{\HMTLang{Composition unitaire avec archive bilatérale}{Unitary composition with bilateral archive}}
\HMTLang{Pour un transport unitaire \(C\) dans un espace de Hilbert \(\mathcal H\), soient}{For a unitary transport \(C\) on a Hilbert space \(\mathcal H\), let}
\[
 T_C=\frac{8I+C}{9},\qquad D_C=\frac{\sqrt8(C-I)}9,
 \qquad R_C=\frac{I+8C}{9},\qquad
 \mathbb U_C=\begin{pmatrix}T_C&D_C\\D_C&R_C\end{pmatrix}.
\]
\begin{theorem}[\HMTLang{Représentation bilatérale et composition complète}{Bilateral representation and complete composition}]
\HMTLang{L'opérateur \(\mathbb U_C\) est unitaire et satisfait}{The operator \(\mathbb U_C\) is unitary and satisfies}
\[
 \mathbb U_{C_2}\mathbb U_{C_1}=\mathbb U_{C_2C_1},\qquad
 T_{C_2}T_{C_1}+D_{C_2}D_{C_1}=\frac{8I+C_2C_1}{9}.
\]
\HMTLang{Pour un état à archive initialement nulle,}{For a state with initially zero archive,}
\[
 \|T_Cv\|^2+\|D_Cv\|^2=\|v\|^2,\qquad
 v=T_C^*T_Cv+D_C^*D_Cv.
\]
\end{theorem}
\begin{proof}
\HMTLang{Les matrices}{The matrices}
\[
 P_0=\frac19\begin{pmatrix}8&-\sqrt8\\-\sqrt8&1\end{pmatrix},
 \qquad P_1=\frac19\begin{pmatrix}1&\sqrt8\\\sqrt8&8\end{pmatrix}
\]
\HMTLang{sont des projecteurs orthogonaux complémentaires : la multiplication vérifie \(P_i^2=P_i\), \(P_0P_1=0\), \(P_0+P_1=I_2\). Puisque \(\mathbb U_C=P_0\otimes I+P_1\otimes C\), son produit adjoint est l'identité et la composition multiplie uniquement les facteurs \(C\). La première composante du produit donne la formule bilatérale. Appliquer l'unitarité à \((v,0)\) prouve l'égalité des normes ; son bloc supérieur gauche est \(T_C^*T_C+D_C^*D_C=I\), qui donne l'application de récupération.}{are complementary orthogonal projectors: multiplication verifies \(P_i^2=P_i\), \(P_0P_1=0\), and \(P_0+P_1=I_2\). Since \(\mathbb U_C=P_0\otimes I+P_1\otimes C\), its adjoint product is the identity, and composition multiplies only the factors \(C\). The first component of the product gives the bilateral formula. Applying unitarity to \((v,0)\) proves norm conservation; its upper-left block is \(T_C^*T_C+D_C^*D_C=I\), which gives the recovery map.}
\end{proof}
\HMTLang{L’ordre de composition est distribué de manière exacte :}{L’ordre de composition est distribué exactement :}
\[
 [T_{C_2},T_{C_1}]=\frac{[C_2,C_1]}{81},\qquad
 [D_{C_2},D_{C_1}]=\frac{8[C_2,C_1]}{81}.
\]
\HMTLang{En additionnant les deux contributions, on obtient \([C_2,C_1]/9\). Les facteurs \(1\) et \(8\) comparent ces termes opératoriels ; leurs interprétations énergétiques dépendent de la forme et du lecteur. L’archive conserve la contribution qu’une composition de contractions visibles omettrait à elle seule.}{L’addition des deux contributions donne \([C_2,C_1]/9\). Les facteurs \(1\) et \(8\) comparent ces termes opératoriels ; leurs interprétations énergétiques dépendent de la forme et du lecteur. L’archive conserve la contribution que la seule composition des contractions visibles omettrait.}

\subsection{\HMTLang{Plan de quart de tour dans le registre dodécaphasique}{Plan de quart de tour dans le registre dodécaphasique}}
\HMTLang{Soit \(C_{12}e_j=e_{j+1}\), avec des indices modulo \(12\), la permutation du calendrier. Le projecteur sur la période quatre et sa composante antipériodique sont}{Soit \(C_{12}e_j=e_{j+1}\), avec des indices modulo \(12\), la permutation du calendrier. Le projecteur sur la période quatre et sa composante antipériodique sont}
\[
 P_4=\frac{I+C_{12}^4+C_{12}^8}{3},\qquad
 Q=P_4\frac{I-C_{12}^6}{2}.
\]
\HMTLang{Définissons}{Define}
\[
 u_0=(1,0,-1,0,1,0,-1,0,1,0,-1,0)^{\mathsf T},
 \quad v_0=C_{12}u_0,\quad W=\frac1{\sqrt6}(u_0,v_0).
\]
\begin{proposition}[\HMTLang{Réalisation exacte du quart de tour}{Réalisation exacte du quart de tour}]
\[
 W^{\mathsf T}W=I_2,\quad WW^{\mathsf T}=Q,\quad
 C_{12}W=WJ_0,\quad C_{12}^3W=-WJ_0,\qquad
 J_0=\begin{pmatrix}0&-1\\1&0\end{pmatrix}.
\]
\end{proposition}
\begin{proof}
\HMTLang{Les vecteurs \(u_0,v_0\) ont six entrées de module un et des supports disjoints, de sorte qu'ils sont orthogonaux et ont pour norme \(\sqrt6\). En outre, \(C_{12}v_0=-u_0\). L'image de \(Q\) contient précisément les vecteurs de période quatre \((a,b,-a,-b)\) répétés trois fois ; elle est donc de dimension deux et coïncide avec l'image de \(W\). Les identités restantes se déduisent par application à ses deux colonnes.}{The vectors \(u_0,v_0\) each have six entries of modulus one and disjoint supports, so they are orthogonal with norm \(\sqrt6\). Moreover, \(C_{12}v_0=-u_0\). The range of \(Q\) consists precisely of period-four vectors \((a,b,-a,-b)\) repeated three times; it is therefore two-dimensional and equals the image of \(W\). The remaining identities follow by applying the operators to both columns.}
\end{proof}
% Fragmento ES para inserción, no documento autónomo.
% Destino: XII, ampliacion_20260922.tex, después de la prueba de
% «Realización exacta del cuarto de giro» y antes de introducir S_zeta.
% Propietarios y comprobaciones: LEER_PRIMERO.md.
\subsubsection{Restitution du calendrier et retenue ternaire}
\label{ins29:xii:acarreo}

L’extraction du plan de quart de tour admet une description exacte
de l’information temporelle conservée. Nous recevons le calendrier de
\(108\) positions et l’extension dodécaphasique sur la phase nonadique
construits précédemment. Pour cette description, nous utilisons des représentants
indexés à partir de zéro :
\[
 \theta=r+9\beta,\qquad 0\leq r<9,\quad 0\leq\beta<12.
\]
La position que le registre énumère de \(1\) à \(9\) est ici \(r+1\).
La mise à jour enrichie précède ce calendrier fini et conserve,
en outre, la route, les incidences et la mémoire accumulée. Dans le plan déjà
construit, \(C_{12}^{\beta}W=WJ_0^{\beta}\) ; sa lecture dépend donc
de \(b=\beta\bmod4\).

\begin{proposition}[Fibre ternaire de la lecture de quart de tour]
\label{ins29:xii:fibra}
Le couple formé par la phase nonadique \(r\) et la phase planaire \(b\) détermine
un quotient de \(36\) positions du calendrier. Sa loi de composition est
\[
 (r,b)(r',b')=
 \left(r+r'\bmod9,\;
 b+b'+\left\lfloor\frac{r+r'}9\right\rfloor\bmod4\right).
\]
Chaque lecture possède trois préimages dans \(\mathbb Z/108\mathbb Z\).
La composante supplémentaire \(k=\beta\bmod3\) restitue la coordonnée
dodécaphasique au moyen de
\[
 \beta=9b+4k\pmod{12}.
\]
\end{proposition}
\begin{proof}
La somme de deux représentants donne
\(\theta+\theta'=r+r'+9(\beta+\beta')\). La division euclidienne de
\(r+r'\) par \(9\) produit la retenue de la formule. La lecture
\((r,b)\) équivaut au résidu \(\theta\bmod36\), puisque
\(\theta\equiv r+9b\pmod{36}\). Son noyau est
\(\{0,36,72\}\), de sorte que chaque fibre a trois éléments.
Enfin, \(9b+4k\) est congru à \(b\) modulo \(4\) et à
\(k\) modulo \(3\) ; ces deux résidus déterminent un unique élément
modulo \(12\).
\end{proof}

\begin{proposition}[Localisation de l’extension non scindée]
\label{ins29:xii:extension}
La décomposition
\[
 \theta\longmapsto(j,z)=(\theta\bmod4,\theta\bmod27),
 \qquad \theta=81j+28z\pmod{108},
\]
identifie le calendrier à
\(\mathbb Z/4\mathbb Z\times\mathbb Z/27\mathbb Z\).
La projection nonadique est \((j,z)\mapsto z\bmod9\), et l’inclusion
dodécaphasique \(\beta\mapsto9\beta\) prend la forme
\[
 \beta\longmapsto(\beta\bmod4,9\beta\bmod27).
\]
L’obstruction à scinder l’extension complète se concentre dans
\[
 0\longrightarrow\mathbb Z/3\mathbb Z
 \longrightarrow\mathbb Z/27\mathbb Z
 \longrightarrow\mathbb Z/9\mathbb Z\longrightarrow0.
\]
\end{proposition}
\begin{proof}
Les coefficients \(81\) et \(28\) ont les résidus \((1,0)\) et \((0,1)\)
modulo \((4,27)\), respectivement. Cela vérifie l’inverse et les deux
applications indiquées. La suite ternaire ne se scinde pas : un produit
\(\mathbb Z/3\mathbb Z\times\mathbb Z/9\mathbb Z\) a pour exposant
\(9\), tandis que \(\mathbb Z/27\mathbb Z\) a pour exposant \(27\).
En revanche, la projection du quotient de \(36\) positions sur
\(\mathbb Z/9\mathbb Z\) admet la section homomorphe
\(n\mapsto28n\pmod{36}\). La réduction élimine donc précisément
cette obstruction ternaire.
\end{proof}

La distinction s’observe déjà dans \(\theta=0,36,72\) : les trois positions
donnent \(r=b=0\), tandis que \(z\) vaut \(0,9,18\), respectivement.
La restitution du calendrier requiert de conserver la fibre que le plan
identifie. En outre, le résidu quaternaire global satisfait
\(j=r+b\pmod4\) ; la phase planaire \(b\) est transportée avec la phase
nonadique et sa retenue.

La portée du quart de tour se trouve ainsi localisée : il organise une
représentation orientée du registre et, accompagné du résidu nonadique,
retrouve exactement son quotient de \(36\) positions. La coordonnée
ternaire supplémentaire restitue les \(108\) positions. La mémoire de la
dynamique enrichie conserve à son tour le prolongement entre les retours
complets du calendrier. La périodicité d’une représentation et la
persistance de la trajectoire enregistrée appartiennent donc à des lectures distinctes
du même transport.

\HMTLang{Le plan est extrait de la permutation déjà définie. Sa forme d’action est reçue du couple angulaire engendré, avec \(A\) et \(C^*\) exprimés dans la même unité, \(|C^*/A|<1\) et \(\theta\in\mathbb R\) :}{The plane is extracted from the already defined permutation. Its action form is received from the generated angular pair, with \(A\) and \(C^*\) expressed in the same unit, \(|C^*/A|<1\), and \(\theta\in\mathbb R\):}
\[
 \xi=C^*/A,\quad \zeta=\operatorname{artanh}\xi,\quad
 S_\zeta=\operatorname{diag}(e^{\zeta/2},e^{-\zeta/2}),\quad
 U_\zeta(\theta)=S_\zeta e^{-\theta J_0}S_\zeta^{-1}
 =\begin{pmatrix}\cos\theta&e^\zeta\sin\theta\\
 -e^{-\zeta}\sin\theta&\cos\theta\end{pmatrix}.
\]
\HMTLang{Chaque \(U_\zeta\) conserve \(g_\zeta=S_\zeta^{-\mathsf T}S_\zeta^{-1}\) et a pour déterminant un. Sa forme conservée est spécifiée avant de comparer des transports d’anisotropies différentes.}{Chaque \(U_\zeta\) préserve \(g_\zeta=S_\zeta^{-\mathsf T}S_\zeta^{-1}\) et a pour déterminant un. Sa forme conservée est spécifiée avant de comparer des transports d’anisotropies différentes.}

\HMTLang{Les produits d’opérateurs agissent de droite à gauche.}{Les produits d’opérateurs agissent de droite à gauche.}

\subsection{\HMTLang{Commutateur relatif et restitution de l’action}{Commutateur relatif et récupération de l’action}}
\begin{theorem}[\HMTLang{Défaut de commutation des formes conjuguées}{Défaut de commutation des formes conjuguées}]
\HMTLang{Soient \(U=U_{\zeta_s}(\theta)\), \(V=U_{\zeta_r}(\phi)\), \(H=UVU^{-1}V^{-1}\). Avec \(\delta=\zeta_s-\zeta_r\) et \(\chi=\sinh\delta\sin\theta\sin\phi\),}{Soient \(U=U_{\zeta_s}(\theta)\), \(V=U_{\zeta_r}(\phi)\) et \(H=UVU^{-1}V^{-1}\). Avec \(\delta=\zeta_s-\zeta_r\) et \(\chi=\sinh\delta\sin\theta\sin\phi\),}
\[
 \det H=1,\qquad \operatorname{tr}H=2+4\chi^2,\qquad
 \lambda_\pm(H)=(\sqrt{1+\chi^2}\pm|\chi|)^2.
\]
\HMTLang{De plus, \(H=I\) exactement lorsque \(\chi=0\).}{De plus, \(H=I\) précisément lorsque \(\chi=0\).}
\end{theorem}
\begin{proof}
\HMTLang{La multiplication des matrices donne \(\operatorname{tr}(UV)=2(\cos\theta\cos\phi-\cosh\delta\sin\theta\sin\phi)\). Pour les matrices de déterminant un, Cayley--Hamilton donne \(A^{-1}=(\operatorname{tr}A)I-A\). Substituer cette expression aux deux inverses produit}{La multiplication donne \(\operatorname{tr}(UV)=2(\cos\theta\cos\phi-\cosh\delta\sin\theta\sin\phi)\). Pour les matrices de déterminant un, Cayley--Hamilton donne \(A^{-1}=(\operatorname{tr}A)I-A\). Substituer cette expression aux deux inverses donne}
\[
 \operatorname{tr}(UVU^{-1}V^{-1})
 =a^2+b^2+c^2-abc-2,
 \quad a=\operatorname{tr}U,\ b=\operatorname{tr}V,\ c=\operatorname{tr}(UV).
\]
\HMTLang{La substitution trigonométrique réduit l’expression à \(2+4\chi^2\). Résoudre \(\lambda^2-(2+4\chi^2)\lambda+1=0\) donne les valeurs propres. Si \(\chi=0\), soit les formes coïncident, soit l’un des transports est \(\pm I\) ; les deux cas commutent. L’implication réciproque se déduit de la trace.}{La substitution trigonométrique réduit cette expression à \(2+4\chi^2\). Résoudre \(\lambda^2-(2+4\chi^2)\lambda+1=0\) donne les valeurs propres. Si \(\chi=0\), soit les formes coïncident, soit un transport est \(\pm I\) ; les deux cas commutent. La réciproque découle de la trace.}
\end{proof}
\HMTLang{Pour les quarts de tour conjugués \(U_+=U_\zeta(\pi/2)\), \(U_-=U_{-\zeta}(\pi/2)\), les inverses sont \(-U_\pm\). Par multiplication,}{For conjugate quarter turns \(U_+=U_\zeta(\pi/2)\), \(U_-=U_{-\zeta}(\pi/2)\), the inverses are \(-U_\pm\). Multiplication gives}
\begin{equation}
 H=(U_+U_-)^2=\operatorname{diag}(e^{4\zeta},e^{-4\zeta}),\qquad
 \operatorname{tr}H-2=\frac{16\xi^2}{(1-\xi^2)^2}.
 \label{delta22:xi:lazo}
\end{equation}
\HMTLang{Avec les représentations \(A=1000\alpha\), \(C^*=2(\eta_{\rm ret}+\alpha)\), toutes deux en degrés, on obtient}{With the outputs \(A=1000\alpha\), \(C^*=2(\eta_{\rm ret}+\alpha)\), both in degrees, this gives}
\begin{equation}
 \rho_{\rm lazo}=e^{4\zeta}
 =\left(\frac{501\alpha+\eta_{\rm ret}}{499\alpha-\eta_{\rm ret}}\right)^2,
 \qquad 0<\eta_{\rm ret}<499\alpha.
 \label{delta22:xi:rho}
\end{equation}
\HMTLang{Son inverse, en conservant l'axe orienté, est}{Its inverse, retaining the oriented axis, is}
\[
 \eta_{\rm ret}=\alpha\left[
 500\frac{\sqrt{\rho_{\rm lazo}}-1}{\sqrt{\rho_{\rm lazo}}+1}-1\right],
 \qquad \hbar_{\rm ret}=s_0\eta_{\rm ret}.
\]
\HMTLang{La preuve utilise \(e^{2\zeta}=(1+\xi)/(1-\xi)\) et isole \(\xi\). On conserve \(s_0=10^{-34}\mathcal U_S\), préalablement construit. Sous conjugaison, l’étiquette de l’axe est transportée ; la seule trace conserve \(|\zeta|\). Le rapport \(\rho_{\rm lazo}\) désigne cette réponse relative, tandis que \(\pi/\varphi^2\) conserve la définition distincte de rapport régional marqué.}{La preuve utilise \(e^{2\zeta}=(1+\xi)/(1-\xi)\) et résout pour \(\xi\). On retient \(s_0=10^{-34}\mathcal U_S\), préalablement construit. Sous conjugaison, l’étiquette de l’axe est transportée ; la seule trace retient \(|\zeta|\). Le rapport \(\rho_{\rm lazo}\) désigne cette réponse relative, tandis que \(\pi/\varphi^2\) conserve sa définition distincte de rapport régional marqué.}

\HMTLang{La branche positive correspond exactement à \(\rho_{\rm lazo}>(501/499)^2\). Lors de l'inversion de l'orientation, \(\rho_{\rm lazo}\mapsto\rho_{\rm lazo}^{-1}\), l'étiquette de l'axe est transportée avant d'appliquer l'application de récupération de cette branche.}{The positive branch corresponds exactly to \(\rho_{\rm lazo}>(501/499)^2\). Under orientation reversal, \(\rho_{\rm lazo}\mapsto\rho_{\rm lazo}^{-1}\), the axis label is transported before applying the recovery map of that branch.}

\subsection{\HMTLang{Relèvement dodécaphasique et amplification compatible}{Relèvement dodécaphasique et amplification compatible}}
\HMTLang{Définissons sur les douze coordonnées}{Définissons sur les douze coordonnées}
\[
 \widetilde S_\zeta=I-Q+WS_\zeta W^{\mathsf T},\qquad
 L_\pm=\widetilde S_{\pm\zeta}C_{12}^3\widetilde S_{\pm\zeta}^{-1}.
\]
\begin{proposition}[\HMTLang{Restitution du commutateur dans le registre complet}{Restauration du commutateur dans le registre complet}]
\[
 L_+L_-L_+^{-1}L_-^{-1}=I-Q+WHW^{\mathsf T}.
\]
\end{proposition}
\begin{proof}
\HMTLang{Les opérateurs \(Q\) et \(C_{12}\) commutent. Sur \(\operatorname{im}Q\), l’isométrie \(W\) identifie \(C_{12}^3\) à \(-J_0\), et \(L_\pm\) à \(U_\pm\). Sur le complément, les deux \(L_\pm\) coïncident avec \(C_{12}^3\), de sorte que leur commutateur est l’identité. La somme orthogonale des deux restrictions donne la formule.}{Les opérateurs \(Q\) et \(C_{12}\) commutent. Sur \(\operatorname{im}Q\), l’isométrie \(W\) identifie \(C_{12}^3\) à \(-J_0\), et \(L_\pm\) à \(U_\pm\). Sur le complément, les deux \(L_\pm\) sont égaux à \(C_{12}^3\), de sorte que leur commutateur est l’identité. La somme orthogonale des deux restrictions donne la formule.}
\end{proof}
\HMTLang{Pour l'amplification déclarée \(\mathcal H_m=\mathbb R^{12}\otimes(\mathbb R^9)^{\otimes m}\), soient}{For the specified amplification \(\mathcal H_m=\mathbb R^{12}\otimes(\mathbb R^9)^{\otimes m}\), let}
\[
 \widetilde H=I-Q+WHW^{\mathsf T},\qquad
 H_m=\widetilde H\otimes I_{9^m},\qquad
 J_m v=v\otimes u_9,\quad u_9=(1,\ldots,1)/3.
\]
\HMTLang{Alors \(J_m\) est isométrique et \(H_{m+1}J_m=J_mH_m\). La norme de l'opérateur reste fixe lors d'une amplification par une identité ; l'action compatible s'étend donc à la limite de Hilbert. Les preuves précédentes sont conservées dans ce raffinement tensoriel, avec les projecteurs et le lecteur transportés.}{Then \(J_m\) is isometric and \(H_{m+1}J_m=J_mH_m\). Operator norm remains fixed under amplification by an identity, so the compatible action extends to the Hilbert limit. The preceding proofs are retained in this tensor refinement, with projectors and reader transported.}

\subsection{\HMTLang{Information chronologique et lecture énergétique}{Information chronologique et lecture énergétique}}
\HMTLang{La puissance de la boucle est \(H^n=\operatorname{diag}(\rho_{\rm lazo}^n,\rho_{\rm lazo}^{-n})\). Elle conserve \(QP\) et \(dQ\wedge dP\). Pour \(\rho_{\rm lazo}\ne1\), une préparation ayant une composante \(Q_0\ne0\) et un axe connu permet de retrouver \(n=\log(Q_n/Q_0)/\log\rho_{\rm lazo}\). La préparation et l’échelle du détecteur appartiennent au domaine de ce récupérateur.}{La puissance de la boucle est \(H^n=\operatorname{diag}(\rho_{\rm lazo}^n,\rho_{\rm lazo}^{-n})\). Elle préserve \(QP\) et \(dQ\wedge dP\). Pour \(\rho_{\rm lazo}\ne1\), une préparation avec \(Q_0\ne0\) et un axe connu retrouve \(n=\log(Q_n/Q_0)/\log\rho_{\rm lazo}\). La préparation et l’échelle du détecteur appartiennent au domaine de cette application de récupération.}

\HMTLang{On admet \(n\in\mathbb Z\), y compris l’itération de la boucle inverse. Lorsque \(Q_0=0\) et \(P_0\ne0\), l’autre composante fournit}{On admet \(n\in\mathbb Z\), y compris l’itération de la boucle inverse. Lorsque \(Q_0=0\) et \(P_0\ne0\), l’autre composante donne}
\[
 n=-\frac{\log|P_n/P_0|}{\log\rho_{\rm lazo}}.
\]

\HMTLang{Les deux ordonnancements \(UU^{-1}VV^{-1}\) et \(UVU^{-1}V^{-1}\) contiennent les mêmes quatre multiplicités, mais leurs réponses sont \(I\) et \(H\). L'archive ordonnée conserve une distinction que le comptage commutatif élimine. Pour un ensemble de détecteurs \(\mathcal D\), deux histoires sont opérationnellement équivalentes lorsque toutes leurs lectures dans \(\mathcal D\) coïncident. Les réponses sur une base complète déterminent l'opérateur linéaire ; la récupération d'une histoire enrichie requiert en outre ses autres registres.}{The two orderings \(UU^{-1}VV^{-1}\) and \(UVU^{-1}V^{-1}\) contain the same four multiplicities, but their responses are \(I\) and \(H\). The ordered archive retains a distinction removed by commutative counting. For a detector set \(\mathcal D\), two histories are operationally equivalent when all their readings in \(\mathcal D\) agree. Responses on a complete basis determine the linear operator; recovery of an enriched history additionally requires its remaining records.}

\HMTLang{La forme de référence fixe la signification énergétique. Pour \(g_{\rm ref}=\operatorname{diag}(e^{-\zeta_r},e^{\zeta_r})\), les gains d’énergie quadratique de la boucle diagonale, sur les axes marqués, sont}{La forme de référence fixe la signification énergétique. Pour \(g_{\rm ref}=\operatorname{diag}(e^{-\zeta_r},e^{\zeta_r})\), les gains d’énergie quadratique de la boucle diagonale le long des axes marqués sont}
\[
 G_Q=e^{8\zeta},\qquad G_P=e^{-8\zeta},\qquad
 \mathcal O=\tfrac14\log(G_Q/G_P)=4\zeta.
\]
\HMTLang{Ces gains sont les carrés des dilatations d’état. \(H\) et \(H^{-1}\) ont la même trace, tandis que \(\mathcal O\) conserve le signe orienté. Pour \(\zeta\ne0\), \(U_+\) et \(U_-\) ne possèdent pas de métrique positive commune : s’ils en conservaient tous deux une, leur commutateur la conserverait également et toutes ses valeurs propres seraient de module un, en contradiction avec \eqref{delta22:xi:lazo}. Chaque transport individuel conserve sa propre \(g_{\pm\zeta}\).}{Ces gains sont les carrés des dilatations d’état. \(H\) et \(H^{-1}\) ont la même trace, tandis que \(\mathcal O\) retient le signe orienté. Pour \(\zeta\ne0\), \(U_+\) et \(U_-\) n’ont pas de métrique positive commune : s’ils en préservaient tous deux une, leur commutateur la préserverait également et toutes ses valeurs propres seraient de module un, en contradiction avec \eqref{delta22:xi:lazo}. Chaque transport individuel préserve sa propre \(g_{\pm\zeta}\).}

\HMTLang{L’identité algébrique de composition bilatérale demeure valable pour ces opérateurs. La conservation unitaire conjointe requiert en outre une même forme positive conservée. L’holonomie relative conserve l’aire et ses propres données de restitution.}{L’identité algébrique de composition bilatérale demeure valable pour ces opérateurs. La conservation unitaire conjointe requiert en outre une forme positive commune préservée. L’holonomie relative préserve l’aire et ses propres données de récupération.}

\HMTLang{Un transport passif \(T_n=S_{\zeta_{n+1}}R(-\theta_n)S_{\zeta_n}^{-1}\) se télescope : son produit est \(S_{\zeta_N}R(-\sum\theta_n)S_{\zeta_0}^{-1}\), égal à l’identité lorsque la forme et la phase totale reviennent. Le commutateur relatif compare, en revanche, des opérations effectives sur une référence commune. Attribuer à cette réponse un échange énergétique exige d’inclure le système, la mémoire et le dispositif. La conservation de l’aire et la réversibilité opératorielle sont démontrées indépendamment de cette attribution.}{Un transport passif \(T_n=S_{\zeta_{n+1}}R(-\theta_n)S_{\zeta_n}^{-1}\) se télescope : son produit est \(S_{\zeta_N}R(-\sum\theta_n)S_{\zeta_0}^{-1}\), égal à l’identité lorsque la forme et la phase totale reviennent. Le commutateur relatif compare en revanche des opérations effectives à une référence commune. Attribuer un échange énergétique à cette réponse exige d’inclure le système, la mémoire et le dispositif. La conservation de l’aire et la réversibilité opératorielle sont prouvées indépendamment de cette attribution.}

\HMTLang{Le résultat réunit trois niveaux de récupération : l'archive bilatérale complète la composition visible ; le calendrier dodécaphasique détermine un plan opératoire exact ; le lecteur orienté de la boucle récupère le contraste angulaire et l'action avec leurs antécédents. Cette articulation prolonge la fonction de couplage de \(K\) sans remplacer sa sélection par une inversion ultérieure. La distinction entre phase, ordre de composition, forme conservée et mémoire récupérable est exprimée au moyen d'opérateurs et d'applications de récupération explicites.}{The result assembles three levels of recovery: the bilateral archive completes visible composition; the dodecaphasic calendar determines an exact operational plane; and the oriented loop reader recovers angular contrast and action with their antecedents. This articulation extends the coupling role of \(K\) without replacing its selection by a later inversion. The distinction between phase, composition order, conserved form, and recoverable memory is expressed by explicit operators and recovery maps.}
\par\endgroup
